\chapter{HISTORICAL NOTE}

The history of Non-Euclidean Geometry has been so well and so often
written that we will give only a brief outline.\label{3hypos}

There is one axiom of Euclid that is somewhat complicated in its
expression and does not seem to be, like the rest, a simple
elementary fact. It is this:\footnote{See article on the axioms of
Euclid by Paul Tannery. \emph{Bulletin des Sciences Mathématiques},
1884.}

\begin{quotation}\emph{If two lines are cut by a third, and the sum of the
interior angles on the same side of the cutting line is less than
two right angles, the lines will meet on that side when sufficiently
produced.}\end{quotation}

Attempts were made by many mathematicians, notably by Legendre, to
give a proof of this proposition; that is, to show that it is a
necessary consequence of the simpler axioms preceding it. Legendre
proved that the sum of the angles of a triangle can never exceed two
right angles, and that if there is a single triangle in which this
sum is equal to two right angles, the same is true of all triangles.
This was, of course, on the supposition that a line is of infinite
length. He could not, however, prove that there exists a triangle
the sum of whose angles is two right angles.\footnote{See, for
example, the twelfth edition of his \emph{Éléments de Géométrie},
Livre~I, Proposition~XIX, and Note~II. See also a statement by Klein
in an article on the Non-Euclidean Geometry in the second volume of
the first series of the \emph{Bulletin des Sciences Mathématiques}.}

At last some mathematicians began to believe that this statement was
not capable of proof, that an equally consistent Geometry could be
built up if we suppose it not always true, and, finally, that all of
the postulates of Euclid were only hypotheses which our experience
had led us to accept as true, but which could be replaced by
contrary statements in the development of a logical Geometry.

The beginnings of this theory have sometimes been ascribed to Gauss,
but it is known now that a paper was written by
Lambert,\footnote{See \emph{American Mathematical Monthly},
July--August, 1895.} in 1766, in which he maintains that the
parallel axiom needs proof, and gives some of the characteristics of
Geometries in which this axiom does not hold. Even as long ago as
1733 a book was published by an Italian, Saccheri, in which he gives
a complete system of Non-Euclidean Geometry, and then saves himself
and his book by asserting dogmatically that these other hypotheses
are false. It is his method of treatment that has been taken as the
basis of the first chapter of this book.\footnote{The translation of
Saccheri by Halsted has been appearing in the \emph{American
Mathematical Monthly}.}

Gauss was seeking to prove the axiom of parallels for many years,
and he may have discovered some of the theorems which are
consequences of the denial of this axiom, but he never published
anything on the subject.

Lobachevsky, in Russia, and Johann Bolyai, in Hungary, first
asserted and proved that the axiom of parallels is not necessarily
true. They were entirely independent of each other in their work,
and each is entitled to the full credit of this discovery. Their
results were published about 1830.

It was a long time before these discoveries attracted much notice.
Meanwhile, other lines of investigation were carried on which were
afterwards to throw much light on our subject, not, indeed, as
explanations, but by their striking analogies.

Thus, within a year or two of each other, in the same journal
(Crelle) appeared an article by Lobachevsky giving the results of
his investigations, and a memoir by Minding on surfaces on which he
found that the formulæ of Spherical Trigonometry hold if we put $ia$
for $a$, etc. Yet these two papers had been published thirty years
before their connection was noticed (by Beltrami).\label{histnote}

Again, Cayley, in 1859, in the \emph{Philosophical Transactions},
published his Sixth Memoir on Quantics, in which he developed a
projective theory of measurement and showed how metrical properties
can be treated as projective by considering the anharmonic relations
of any figures with a certain special figure that he called the
absolute. In 1872 Klein took up this theory and showed that it gave
a perfect image of the Non-Euclidean Geometry.

It has also been shown that we can get our Non-Euclidean Geometries
if we think of a unit of measure varying according to a certain law
as it moves about in a plane or in space.

The older workers in these fields discovered only the Geometry in
which the hypothesis of the acute angle is assumed. It did not occur
to them to investigate the assumption that a line is of finite
length. The Elliptic Geometry was left to be discovered by Riemann,
who, in 1854, took up a study of the foundations of Geometry. He
studied it from a very different point of view, an abstract
algebraic point of view, considering not our space and geometrical
figures, except by way of illustration, but a system of variables.
He investigated the question, What is the nature of a function of
these variables which can be called element of length or distance?
and found that in the simplest cases it must be the square root of a
quadratic function of the differentials of the variables whose
coefficients may themselves be functions of the variables. By taking
different forms of the quadratic expressions we get an infinite
number of these different kinds of Geometry, but in most of them we
lose the axiom that bodies may be moved about without changing their
size or shape.

Two more names should be included in this sketch,---Helmholtz and
Clifford. These did much to make the subject popular by articles in
scientific journals. To Clifford we owe the theory of parallels in
elliptic space, as explained on page~\pageref{p68ref}. He showed
that we can have in this Geometry a finite surface on which the
Euclidean Geometry holds true.\footnote{Some of the more interesting
accounts of Non-Euclidean Geometry are: \emph{Encyclopedia
Britannica}, article ``Measurement'', by Sir Robert Ball.
\emph{Revue Générale des Sciences}, 1891, ``Les Géométries
Non-Euclidean,'' by Poincaré. \emph{Bulletin of the American
Mathematical Society}, May and June, 1900, ``Lobachevsky's
\emph{Geometry}'', by Frederick S.~Woods. \emph{Mathematische
Annalen}, Bd.~xlix, p.~149, 1897, and \emph{Bulletin des Sciences
Mathématiques}, 1897, ``Letters of Gauss and Bolyai''; particularly
interesting is one letter in which Gauss gives a formula for the
area of a triangle on the hypothesis that we can draw three mutually
parallel lines enclosing a finite area always the same. The last two
articles refer to the publications of Professors Engel and Stäckel,
which give in German a full history of the theory of parallels and
the writings and lives of Lobachevsky and Bolyai. See also the
translations by Prof. George Bruce Halsted of Lobachevsky and Bolyai
and of an address by Professor Vasiliev.}

The chief lesson of Non-Euclidean Geometry is that the axioms of
Geometry are only deductions from our experience, like the theories
of physical science. For the mathematician, they are hypotheses
whose truth or falsity does not concern him, but only the
philosopher. He may take them in any form he pleases and on them
build his Geometry, and the Geometries so obtained have their
applications in other branches of mathematics.

The ``axiom'', so far as this word is applied to these geometrical
propositions, is not ``self-evident'', and is not necessarily true.
If a certain statement can be proved,---that is, if it is a
necessary consequence of axioms already adopted,---then it should
not be called an axiom. When two or more mutually contradictory
statements are equally consistent with all the axioms that have
already been accepted, then we are at liberty to take either of
them, and the statement which we choose becomes for our Geometry an
axiom. Our Geometry is a study of the consequences of this axiom.

The assumptions which distinguish the three kinds of Geometry that
we have been studying may be expressed in different forms. We may
say that one or two or no parallels can be drawn through a point;
or, that the sum of the angles of a triangle is equal to, less than,
or greater than two right angles; or, that a straight line has two
real points, one real point, or no real point at infinity; or, that
in a plane we can have similar figures or we cannot have similar
figures, and a straight line is of finite or infinite length, etc.
But any of these forms determines the nature of the Geometry, and
the others are deducible from it.

\newpage
